\subsection{Where the Existing Instruments Reach}
\label{sec:12.2.2}
The areas below do not rank or sequence against each other: a system meeting the bar in one and lagging in another is not partway through a queue, because nothing establishes an order for them to arrive in.

\runin{Theory of mind}
The milestone is met when a system's theory of mind holds up under interaction rather than narration. FANToM, built from conversations where speakers join and leave holding different beliefs, found that systems scoring well on passive belief-tracking narratives degrade sharply once the same reasoning runs inside a live, information-asymmetric exchange \autocite{kim2023fantom}. That degradation is what a claimed milestone here has to survive.

\runin{Moral reasoning}
Whether a balance between fast and slow moral judgment survives an actual incentive to cut corners is not a question a static ethics quiz can answer. The MACHIAVELLI benchmark tests it directly, scoring agents across 134 text-based games on both reward and documented violations — deception, harm, power-seeking — and systems optimizing for reward alone drift toward the ends-justify-the-means behavior an antifascist design is meant to rule out \autocite{pan2023machiavelli}. A passing score on a test that applies no pressure is not evidence about a system that will face pressure.

\runin{Robustness}
Here the checking party has to be adversarial rather than merely independent. The narrow slice that can be checked now is whether a trained-in commitment survives someone trying to remove it, measured in red-teaming hours logged and adversarial fine-tuning attempts survived. The milestone is met when a system has been through documented adversarial testing rather than the tests its own designers thought to run, and only once the adversary modeled is a hostile actor deliberately repurposing the system.

\runin{Transparency}
This is the case where the standard is easiest to fake, because a transparency claim is itself a claim about disclosure and can be satisfied in appearance by disclosing something. The milestone counts as met once an audit or review board — the EU AI Act's requirements for high-risk systems are the working instance — has verified it independently, not when a developer's own documentation asserts it.

\runin{Applications}
Perspective API, vTaiwan, the Carnegie Endowment's AI Global Surveillance Index, UNESCO's 2021 Recommendation on the Ethics of Artificial Intelligence: deployments and instruments rather than hypotheticals, and every one of them dual use. What this area needs is a sustained track record, evaluated by a body other than the one that deployed it, for what a system did rather than what it was built to do, over years instead of a launch announcement.

None of the five reaches the floor. Each scores what a system does against conditions its testers chose, and the question — whether the mechanism producing the behavior was tracking the rationale — is the one none of these instruments is built to separate. Robustness comes nearest and stops in a definite place. A commitment that survives an attack has been shown to be durable, and durability is a fact about the constraint's custody rather than about the system holding it. What the standard does not answer is who is qualified to do the checking, and what happens when nobody with the competence to check is also free of a stake in the result.
